\documentclass[conference]{IEEEtran}
\IEEEoverridecommandlockouts
\usepackage{cite}
\usepackage{amsmath,amssymb,amsfonts}
\usepackage{algorithmic}
\usepackage{graphicx}
\usepackage{textcomp}
\usepackage{xcolor}
\usepackage{booktabs}
\usepackage{multirow}
\usepackage{url}
\def\BibTeX{{\rm B\kern-.05em{\sc i\kern-.025em b}\kern-.08em
    T\kern-.1667em\lower.7ex\hbox{E}\kern-.125emX}}
\begin{document}

\title{Phishing Detection with a Multi-Modal Machine Learning System:\\
Integrating URL, Text, Visual, and HTML Signals}

\author{\IEEEauthorblockN{DSA4266 Project Group 5}
\IEEEauthorblockA{\textit{Sense-making Case Analysis: Science and Technology} \\
\textit{National University of Singapore}\\
Singapore}
}

\maketitle

\begin{abstract}
Phishing websites remain one of the most prevalent cyber threats and
traditional blacklist-based defences struggle to keep pace with the
rate at which malicious sites appear. Machine learning offers a more
adaptive alternative, but most existing detectors rely on a single
source of information, typically the URL or the webpage text. In this
work we design, implement, and critically evaluate a multi-modal
phishing detection system that fuses four complementary signals: the
raw URL character sequence together with nine hand-crafted lexical
features, the cleaned visible text of the webpage encoded by
DistilBERT, a screenshot of the rendered page encoded by
EfficientNet-B0, and eight structural features extracted from the raw
HTML. Each modality is projected into a common $256$-dimensional
space and combined through a learned fusion module before a
three-layer MLP bottleneck classifier. On a dataset of 1{,}608
matched webpage records we train and evaluate four unimodal baselines
and the multi-modal fusion model in two regimes: a
single-trial \emph{baseline} using default hyperparameters, and a
20-trial \emph{Optuna-tuned} search that promotes the best unimodal
encoders into the fusion backbone.  The tuned multi-modal model
achieves \textbf{F1 $=0.894$}, \textbf{ROC-AUC $=0.957$},
\textbf{PR-AUC $=0.946$}, and \textbf{accuracy $=0.927$}, outperforming
every unimodal branch --- URL-only (F1 $=0.827$, ROC-AUC $=0.940$),
text-only (F1 $=0.726$, ROC-AUC $=0.875$), visual-only
(F1 $=0.538$, ROC-AUC $=0.738$), and HTML-only (F1 $=0.585$,
ROC-AUC $=0.706$) --- and lifting F1 by $+0.087$ over the
baseline attention fusion ($0.807$).  The tuned fusion also attains
the highest precision ($0.905$) and the fewest false positives (8),
confirming that fusion is genuinely more selective than any single
modality.  A formal leave-one-out ablation on the tuned fusion ranks
modality importance by composite-score drop: URL ($-0.103$) is by far
the most important, followed by HTML ($-0.015$) and text ($-0.003$),
while disabling visual actually \emph{improves} composite score
($+0.017$), showing that the frozen EfficientNet-B0 branch adds more
noise than signal on this corpus.  A score-level SHAP analysis
($R^{2}=0.955$ surrogate fidelity) and a modality-level Shapley
analysis on the projected fusion vectors corroborate the ablation:
URL is the dominant contributor, followed by visual, text, and HTML.
Beyond raw accuracy we perform an exploratory data study, a
temperature-calibration analysis, a threshold sweep, and an
embedding $t$-SNE visualisation. The results reveal that a
carefully balanced fusion architecture (matched classifier capacity
across branches, a large-enough projection dimension to avoid an
information bottleneck, and a composite-score-driven hyperparameter
search) beats the strongest unimodal baseline on every headline
metric, motivating future work on larger corpora, adversarial
robustness, and fine-tuned encoders.
\end{abstract}

\begin{IEEEkeywords}
phishing detection, multi-modal learning, attention fusion,
DistilBERT, EfficientNet, BiLSTM, interpretability, cybersecurity,
class imbalance, threshold calibration
\end{IEEEkeywords}

\section{Introduction}
Phishing is a class of social-engineering attack in which adversaries
impersonate trusted brands through fraudulent URLs, emails, or
websites in order to harvest credentials or financial information.
The Anti-Phishing Working Group has consistently reported several
hundred thousand unique phishing sites every quarter, and the typical
lifespan of an active phishing domain is measured in hours rather
than days. This short window makes reactive defences such as URL
blacklists fundamentally insufficient: a freshly registered site is
usually unknown to any reputation service at the moment a user is
tricked into visiting it.

Machine-learning-based detectors respond to this weakness by
learning generalisable patterns of ``phishiness'' directly from data.
Historically, however, most learned detectors have been built on a
\emph{single} modality --- either a character-level model of the URL
or a text classifier operating on the visible body of the page.
In reality, phishing is a layered form of deception. The URL may
contain suspicious sub-domains or brand look-alikes; the HTML may
include password forms, hidden iframes, and meta-refresh redirects;
the rendered text may use urgency or impersonation cues; and the
screenshot may visually mimic a legitimate login page. An attacker
who obfuscates one of these layers cannot easily obfuscate them all.
This observation motivates multi-modal detection: by combining
evidence from several views of a website, a model should be more
robust and, in principle, more accurate.

Our project has three goals that go beyond reporting a single
accuracy number. First, we build a reproducible multi-modal
detection pipeline that jointly consumes URL, text, screenshot, and
HTML-structural information. Second, we \emph{study} how and why
each modality contributes, through an exploratory data study,
ablation experiments, and fusion-strategy comparisons. Third, we
evaluate the practical trade-off between the additional complexity
of multi-modal learning and the security-critical requirements of
high recall, robustness to missing signals, interpretable
decisions, and reasonable inference cost.

The remainder of the paper is organised as follows. Section II
states the problem precisely. Section III reviews applicable
techniques. Section IV describes the dataset and data study,
Section V details the methodology, and Section VI specifies the
experimental setup and evaluation metrics. Section VII reports the
empirical results, Section VIII discusses them with a focus on
interpretability and training dynamics, and Section IX concludes.

\section{Problem Statement}
\subsection{Task definition}
We address binary classification of webpages as either
\emph{phishing} (label 1) or \emph{legitimate} (label 0). For each
webpage we observe four inputs: (i)~the raw URL string, (ii)~the raw
HTML document, (iii)~the cleaned visible text extracted from the
HTML, and (iv)~a full-page rendered screenshot. Given these four
signals the task is to produce a calibrated probability
$\hat{p}\in[0,1]$ that the page is phishing, together with a
threshold-based decision $\hat{y}\in\{0,1\}$.

\subsection{Why this is hard}
Phishing detection is a genuinely difficult supervised learning
problem for four reasons:

\noindent\textbf{Adversarial distribution shift.} Attackers actively
mutate the features on which detectors rely. Any model that depends
on a single brittle cue is trivially evadable.

\noindent\textbf{Class imbalance and cost asymmetry.} In practice,
the vast majority of webpages a user visits are legitimate. Even on
a relatively balanced research dataset a naive accuracy-maximising
classifier can achieve high accuracy by simply predicting the
majority class. Worse, the costs of the two error types are not
symmetric: a false negative (missed phishing) exposes the user to
credential theft, whereas a false positive (blocked legitimate site)
merely causes inconvenience. Our evaluation must reflect that
asymmetry.

\noindent\textbf{Heterogeneous, multi-modal data.} URL, text,
screenshot, and HTML live in four completely different
representation spaces (character sequences, sub-word tokens, raw
pixels, and numeric structural features). Combining them requires
careful architectural design and calibration.

\noindent\textbf{Limited labelled data.} The matched, multi-modal
corpus available to us contains only 1{,}608 fully labelled webpages,
which strongly constrains the complexity of the models we can train
from scratch.

\subsection{Research questions}
These considerations lead directly to four research questions that
structure the rest of the paper:
\begin{itemize}
  \item \textbf{RQ1}: Which modality provides the strongest standalone
  signal for phishing detection?
  \item \textbf{RQ2}: Does multi-modal fusion improve detection over
  the best single-modality baseline?
  \item \textbf{RQ3}: Which fusion strategy (concatenation, weighted
  sum, attention) gives the best accuracy--complexity trade-off?
  \item \textbf{RQ4}: How much does each modality contribute, as
  measured by leave-one-out ablation, and does the model degrade
  gracefully when a modality is weak or absent?
\end{itemize}

\section{Review of Applicable Techniques}
Given the heterogeneous nature of our inputs, several families of
techniques are applicable. We survey them briefly and justify our
choices.

\subsection{URL-based detection}
\noindent\textbf{Hand-crafted features.} Classical detectors rely on
lexical and host-based features: URL length, number of dots, number
of hyphens, presence of an IP address, HTTPS indicator, sub-domain
count, and entropy of the sub-domain string. These features are
cheap, interpretable, and remain competitive baselines when fed to
random forests or gradient-boosted trees.

\noindent\textbf{Character-level neural models.} URLNet
\cite{urlnet} and its successors replace hand-crafted features with
an end-to-end character-level CNN or LSTM. These models capture
patterns that engineered features miss, such as unusual character
$n$-grams in obfuscated domains, and they generalise better to
previously unseen URLs. We adopt this approach in combination with
the hand-crafted features: the learned encoder captures flexible
patterns while the scalar features provide a strong, interpretable
lexical prior that is particularly useful when training data is
scarce.

\subsection{Content- and HTML-based detection}
\noindent\textbf{Bag-of-words and TF--IDF.} Early text-based
detectors used token counts or TF--IDF vectors over the visible
text. They are fast and interpretable but brittle to vocabulary
shift.

\noindent\textbf{Pre-trained Transformers.} Modern systems fine-tune
BERT \cite{bert} or lighter variants such as DistilBERT
\cite{distilbert}, which capture urgency cues, brand mentions, and
the semantic tone of deceptive text. Given our small dataset we
\emph{freeze} DistilBERT and train only a projection head on top,
which is a form of linear probing and avoids over-fitting the
encoder.

\noindent\textbf{HTML-structural features.} Orthogonal to text
semantics, several studies exploit structural properties of the raw
HTML: the number of \texttt{<form>} and \texttt{<input>} elements,
the presence of password fields, the ratio of external to internal
links, and the use of \texttt{meta-refresh} directives. These
features are cheap, robust to text obfuscation, and highly
interpretable. We use eight such features as a fourth, deliberately
simple modality.

\subsection{Visual detection}
\noindent\textbf{Perceptual hashing and logo matching.} Early work
used SIFT descriptors or perceptual hashes to match screenshots
against a database of known brand logos, but these approaches
generalise poorly beyond their reference set.

\noindent\textbf{Pre-trained CNNs.} Recent work fine-tunes a deep
CNN, such as ResNet or EfficientNet \cite{effnet}, on labelled
screenshots. We use a frozen EfficientNet-B0 backbone because it
has substantially fewer parameters than ResNet-50 (5.3M vs.\ 25.6M)
while remaining competitive on image classification and,
importantly, fits comfortably in 3.7~GB of RAM.

\subsection{Multi-modal fusion}
Multi-modal fusion is well studied outside cybersecurity
\cite{multimodal-survey}. Three canonical strategies exist:
\emph{early fusion}, which concatenates raw features before a
single model; \emph{late fusion}, which averages or votes over
per-modality classifier outputs; and \emph{intermediate fusion},
which combines per-modality embeddings through learnable weights
or attention. Intermediate fusion usually performs best because
it preserves modality-specific structure while still allowing
cross-modal interaction. We implement three intermediate-fusion
variants --- concatenation, learnable weighted sum, and multi-head
attention \cite{attn} --- and compare them empirically.

\subsection{Handling small, imbalanced datasets}
Finally, several families of techniques are directly applicable to
the specific statistical challenges of our data:
class-weighted cross-entropy, weighted random samplers, threshold
tuning, temperature calibration, mixup-style regularisation, data
augmentation, early stopping, dropout, and transfer learning
through frozen pre-trained backbones. We incorporate every one of
these in our pipeline.

\section{Dataset and Data Study}
We conduct the exploratory analysis in
\texttt{notebooks/EDA.ipynb}; all numbers in this section are drawn
directly from that notebook's outputs.

\subsection{Data sources and merging}
We combine three public datasets. (i)~A HuggingFace URL dataset
containing $3{,}000$ URL/label pairs split $50{:}50$ between
legitimate and phishing classes. (ii)~A Kaggle HTML corpus with
$1{,}877$ files organised into \texttt{genuine\_site\_0} ($1{,}321$
files) and \texttt{phishing\_site\_1} ($556$ files). (iii)~A Kaggle
screenshot corpus with $1{,}697$ full-page images
($1{,}147$ genuine, $550$ phishing). Filename suffixes in the
Kaggle folders correspond to HuggingFace row indices, which lets us
inner-join the three datasets per sample. The HTML-screenshot
filename overlap is $1{,}684$ of $1{,}877$ (EDA.ipynb cell 29), and
after requiring \emph{all} three modalities for every record the
matched corpus contains $1{,}608$ rows. The HuggingFace train/val/test
split is preserved, yielding final splits of $1{,}128/234/246$
samples.

\subsection{Class balance}
The three raw datasets have different balances: URL is $50{:}50$ by
construction, while the Kaggle HTML and screenshot corpora are
approximately $70{:}30$ in favour of legitimate. The inner join
therefore inherits the Kaggle imbalance: as shown in
Fig.~\ref{fig:classdist}, the matched corpus contains $1{,}103$
legitimate and $505$ phishing pages, approximately $2.2{:}1$. The
test split contains $160$ legitimate and $86$ phishing pages. This
moderate imbalance is enough to skew accuracy-based evaluation: a
trivial constant ``always legitimate'' classifier would achieve
$160/246\approx65\%$ accuracy on the test set while catching zero
phishing pages. This single observation is what forces us to
privilege F1-score, recall, PR-AUC, and threshold calibration over
raw accuracy in Sections VI and VII.

\subsection{Data quality issues}
The raw dataset is not pristine. Three concrete issues surfaced
during preprocessing:

\noindent\textbf{Filename collisions.} A non-trivial number of
Kaggle files use Windows-incompatible filename characters, and
silently fail to read under the native Windows file-system. Our
loader catches these cases with a \texttt{try/except OSError} and
reports the number of unreadable files as a warning.

\noindent\textbf{Missing modalities.} Not every URL has a matched
HTML file or a matched screenshot. We use an \emph{inner join} so
that the final dataset contains only fully-matched examples; the
alternative (imputing missing modalities with zeros) was tested
and found to hurt performance.

\noindent\textbf{Parser failures.} A small fraction of HTML files
contain malformed markup that BeautifulSoup cannot parse. For those
cases we fall back to a regex-based tag stripper so that the text
pipeline always returns a string.

\subsection{Exploratory analysis}
We conducted an exploratory analysis before any modelling. Key
observations, visualised in Figs.~\ref{fig:classdist}
and~\ref{fig:featurecorr}, include:

\begin{figure}[t]
\centering
\includegraphics[width=0.95\columnwidth]{outputs/evaluation/class_distribution.png}
\caption{Class distribution of the matched $1{,}608$-record corpus.
Left: overall split ($1{,}103$ legitimate vs.\ $505$ phishing).
Right: per-split counts. The $\approx 2.2{:}1$ imbalance is inherited
from the Kaggle HTML/screenshot corpora and motivates class-weighted
loss and threshold calibration.}
\label{fig:classdist}
\end{figure}

\noindent\textbf{URL length and structural density.} Phishing URLs
are noticeably longer than legitimate ones: on the raw URL dataset,
mean URL length is $60.3$ characters for phishing vs.\ $44.6$ for
legitimate (EDA.ipynb cell 6). The Pearson correlation with the
phishing label is strongest for \texttt{num\_dots} ($+0.267$),
followed by \texttt{num\_digits} ($+0.181$), \texttt{has\_ip}
($+0.179$), \texttt{subdomain\_count} ($+0.164$),
\texttt{domain\_length} ($+0.161$), \texttt{url\_length} ($+0.159$),
and \texttt{num\_slashes} ($+0.138$). Individually modest, but
jointly informative.

\noindent\textbf{Hyphens.} \texttt{num\_hyphens} has a \emph{negative}
correlation with phishing ($-0.143$), because legitimate SEO-friendly
URLs use hyphens as word separators more often than phishing URLs do.

\noindent\textbf{HTTPS.} The \texttt{has\_https} feature is constant
zero across our URL corpus (the HuggingFace dump stripped the
\texttt{https://} prefix), so its correlation is undefined (NaN).
The feature is therefore non-informative on this dataset and appears
as a blank row in Fig.~\ref{fig:featurecorr}.

\noindent\textbf{HTML structure (counter-intuitive).} On the HTML
corpus (EDA.ipynb cell 17), \emph{every} structural feature is
\emph{negatively} correlated with the phishing label:
\texttt{has\_login\_text} ($r=-0.260$), \texttt{content\_length}
($-0.255$), \texttt{num\_tags} ($-0.204$), \texttt{num\_scripts}
($-0.170$), \texttt{num\_external\_links} ($-0.169$),
\texttt{num\_forms} ($-0.162$), \texttt{num\_images} ($-0.158$),
\texttt{num\_inputs} ($-0.137$), \texttt{num\_iframes} ($-0.122$),
and \texttt{has\_password\_field} ($-0.058$). The legitimate half of
the HTML corpus is dominated by full-featured content sites with
richer HTML, while many phishing pages are short-lived stubs that
fail to render a full login form. Crucially, the sign of these
correlations does \emph{not} mean the features are useless --- it
means the model must \emph{invert} them, which a neural classifier
does easily but a naive rule-based detector would get wrong.

\noindent\textbf{Text length.} Phishing HTML pages have
\emph{less} visible text than legitimate pages
(\texttt{content\_length} correlation $=-0.255$). This is consistent
with the observation that many phishing kits are thin wrappers
around a login form.

\begin{figure}[t]
\centering
\includegraphics[width=0.95\columnwidth]{outputs/evaluation/feature_correlation.png}
\caption{Feature correlation heatmap computed on the matched training
corpus (post-inner-join, $1{,}608$ samples). Numbers differ slightly
from the EDA-notebook correlations (which are computed on the raw
$3{,}000$-URL and $1{,}877$-HTML corpora separately) because the
inner join changes the label distribution. The same qualitative
picture holds: URL structural density correlates positively with
phishing, HTML structural features correlate negatively, and
\texttt{has\_https} is non-informative on this corpus.}
\label{fig:featurecorr}
\end{figure}

\subsection{Preprocessing pipelines}
Based on the data study, we defined four preprocessing pipelines.
URLs are lower-cased, truncated to 200 characters, and character-
encoded through an 85-entry vocabulary; shorter URLs are
right-padded with a reserved token. In parallel, nine normalised
hand-crafted URL features are computed (length, counts of dots,
hyphens, slashes, digits and special characters, IPv4 indicator,
HTTPS indicator, and sub-domain count). Raw HTML is parsed with
BeautifulSoup (\texttt{lxml}), with \texttt{<script>},
\texttt{<style>}, \texttt{<meta>}, \texttt{<noscript>}, and
\texttt{<head>} tags removed; the remaining visible text is
whitespace-normalised and tokenised by the DistilBERT uncased
tokenizer up to 256 sub-word tokens. Screenshots are resized to
$224\times224$ and normalised with ImageNet statistics; during
training we apply stochastic augmentation (random-resized crop in
$[0.85,1.0]$, affine rotation up to $10^{\circ}$, small translation
and zoom, horizontal flip $p{=}0.5$, and mild colour jitter).
Finally, eight structural HTML features are extracted and clipped
to $[0,1]$.

\section{Methodology and Approach}
\subsection{Overall architecture}
Each of the four modalities is first encoded by a dedicated unimodal
feature extractor. The resulting embeddings are passed through
small MLP projection heads that map them into a common $d=256$
space. A fusion module combines the four projected vectors into a
single representation, which is then classified by a three-layer
MLP bottleneck with dropout. We deliberately match the classifier
capacity across the unimodal and multimodal branches so that any
performance gap is attributable to the fusion mechanism rather than
to an accidental capacity asymmetry.

\begin{figure}[t]
\centering
\includegraphics[width=0.95\columnwidth]{outputs/evaluation/embedding_tsne.png}
\caption{$t$-SNE projection of the combined text and visual
embedding space. Phishing and legitimate pages partially separate
into distinct regions, but there is substantial overlap --- which
is exactly why a learned fusion classifier, rather than simple
geometric separability, is needed.}
\label{fig:tsne}
\end{figure}

\subsection{URL encoder: what we tried}
We explored three URL encoder families, all implemented behind a
common interface and selected via the
\texttt{url.model\_type} config flag:
\begin{itemize}
  \item A 1-D character CNN with multiple filter widths.
  \item A two-layer bidirectional LSTM.
  \item A small Transformer with learned positional embeddings.
\end{itemize}
The BiLSTM variant gave the best validation F1 and the most stable
training, and was retained as the default. It uses a 64-dim
character embedding, 128 hidden units per direction, followed by a
linear projection to 256 dimensions and a LayerNorm. We then
concatenate the nine hand-crafted URL features with this learned
256-dim vector before the projection head. This hybrid design
combines the flexibility of a learned encoder with the sample-
efficiency of interpretable features.

\subsection{Text encoder}
The text encoder is a frozen DistilBERT-base-uncased model. We use
mean-pooling over the last hidden state rather than the
\texttt{[CLS]} token, because mean-pooling was empirically more
stable on the short webpage excerpts in our corpus. The 768-dim
output is passed through a trainable projection head. Freezing
DistilBERT dramatically reduces training cost and eliminates the
risk of catastrophic forgetting on our small dataset while
preserving the linguistic knowledge learned during pre-training.

\subsection{Visual encoder}
The visual encoder is a frozen EfficientNet-B0 \cite{effnet}
backbone pre-trained on ImageNet, with its classification layer
removed. The 1280-dim image embedding is passed through a
trainable projection head. We also experimented with ResNet-50 as
a visual backbone but EfficientNet-B0 gave comparable accuracy at
roughly one-fifth the parameter count, and was retained for
efficiency.

\subsection{HTML-structural encoder}
The eight hand-crafted HTML features are fed directly into a small
MLP projection head. No deep encoder is used; the feature set is
already low-dimensional and interpretable, and an additional deep
encoder would simply risk over-fitting.

\subsection{Projection heads and fusion}
Each modality-specific embedding is projected into a common
$d=256$-dim space through a block
$\text{Linear}\rightarrow\text{ReLU}\rightarrow\text{Dropout}\rightarrow
\text{LayerNorm}$. The LayerNorm stabilises training and ensures
the four vectors have comparable scale, which is critical for
weighted and attention-based fusion. An earlier configuration used
$d=128$, but we found that compressing the $265$-dimensional URL
representation (256 learned dims + 9 hand-crafted features) down to
$128$ acted as an information bottleneck that hurt fusion
performance; enlarging the projection to $256$ removes this effect.

We implement three fusion strategies:

\noindent\textbf{Concatenation (baseline).} The four projected
vectors are concatenated into a $1{,}024$-dim vector ($4\times 256$)
and passed to the classifier. This is the simplest form of
intermediate fusion.

\noindent\textbf{Weighted sum.} A learnable vector
$\mathbf{w}\in\mathbb{R}^{4}$ is passed through a softmax and used
to weight the four embeddings:
$\mathbf{z}=\sum_{m}\text{softmax}(\mathbf{w})_{m}\,\mathbf{e}_{m}$.
The fused vector has the same size as any single modality. This
scheme also exposes the learned per-modality weights for
interpretability.

\noindent\textbf{Attention fusion (default).} The four projected
vectors are stacked as a length-four sequence and passed through a
four-head multi-head self-attention layer followed by a residual
connection and LayerNorm. The output is mean-pooled across
modalities. This allows the model to learn cross-modal
interactions rather than treating each modality independently.

\subsection{Classifier head}
The fused representation is passed through a three-layer bottleneck
MLP ($256\rightarrow256\rightarrow128\rightarrow2$) with ReLU
activation and dropout $0.3$. This matches the three-layer bottleneck
used by every unimodal baseline, so that fusion performance is not
inflated (or deflated) by an accidental capacity difference. Training
uses cross-entropy; for evaluation we apply a softmax and treat
$P(\text{phishing})$ as the decision score.

\subsection{How we addressed typical data-science problems}
We deliberately went beyond ``plug in a model'' and applied several
mitigations for the statistical challenges identified in Section IV.

\noindent\textbf{Class imbalance.} We use class-weighted cross-
entropy with weights inversely proportional to class frequencies,
and we calibrate the decision threshold post-hoc on the validation
set (Section VI).

\noindent\textbf{Small-data overfitting.} We freeze the two large
pre-trained backbones (DistilBERT, EfficientNet-B0), train only
lightweight projection heads and a classifier, apply dropout of
$0.3$ and LayerNorm throughout the fusion stack, and use early
stopping with patience 10 on validation loss.

\noindent\textbf{Regularisation beyond dropout.} Image-space
augmentation (random-resized crop, affine rotation, horizontal flip,
colour jitter) is applied to all visual inputs during training.
We implemented manifold mixup inside the fusion stack but disabled
it for the final reported runs after noticing that applying mixup
only to the multimodal training path was giving the multimodal
model an unfair regularisation advantage over the unimodal
baselines; removing it ensures that all models are trained under
identical objectives.

\noindent\textbf{Probability calibration.} Cross-entropy losses
trained with class weights tend to produce over-confident logits.
After training, we perform \emph{temperature scaling} on the
validation set --- a single scalar $T$ is optimised to minimise
validation NLL --- so that reported probabilities are well
calibrated. This is particularly important because our downstream
threshold sweep relies on meaningful probability scores.

\noindent\textbf{Hyperparameter search.} We wrap the training
pipeline in an Optuna-based search driven by a composite metric
$s = 0.5\cdot\text{F1} + 0.25\cdot\text{ROC-AUC} +
0.25\cdot\text{C-index}$ that rewards both classification quality
and probability ranking. Trial pruning is enabled to abort
unpromising runs early. The top-$k$ unimodal configurations are
promoted to the fusion search, so that the final fusion model is
built on top of already-tuned encoders.

\noindent\textbf{Reproducibility.} We fix the seed at 42 across
NumPy, PyTorch, and Python's \texttt{random}, and log every
configuration and result under \texttt{outputs/run.log}. All
experiments can be reproduced end-to-end from
\texttt{scripts/train\_multimodal.py} and the config in
\texttt{configs/config.yaml}.

\noindent\textbf{Interpretability.} We design the system so that
multiple sources of explanation are available: (i)~the nine
hand-crafted URL features and eight HTML features are inherently
interpretable; (ii)~the weighted-fusion variant exposes per-modality
softmax weights; (iii)~the attention fusion variant exposes a
$4\times 4$ attention weight matrix; (iv)~the fused embedding is
projected to 2-D with $t$-SNE; (v)~the top false positives and
false negatives are examined manually (Section VII-F).

\section{Experimental Setup and Evaluation Metrics}
\subsection{Training protocol}
All models are trained with AdamW, learning rate
$5\!\times\!10^{-4}$, weight decay $5\!\times\!10^{-4}$, batch size
16, and a ReduceLROnPlateau scheduler with minimum LR $10^{-6}$.
We use class-weighted cross-entropy, gradient clipping at 1.0, and
early stopping with patience 10 on validation loss. Mixed precision
is enabled when a GPU is available. Every experiment uses the same
seed.

\subsection{Decision threshold calibration}
Because missing a phishing site is far more costly than incorrectly
flagging a legitimate one, the default $0.5$ decision threshold is
not necessarily optimal. After training we sweep the threshold
over $[0.05,0.95]$ on the validation set and select the value that
maximises validation F1-score. For the \emph{baseline} configuration
the selected thresholds are $T^{*}=0.61$ (multi-modal attention
fusion), $0.36$ (URL-only), $0.46$ (text-only), $0.59$ (visual-only),
and $0.58$ (HTML-only). After Optuna tuning the thresholds shift to
$0.43$ (multi-modal weighted fusion), $0.28$ (URL-only), $0.42$
(text-only), $0.58$ (visual-only), and $0.46$ (HTML-only) --- the
tuned fusion threshold moves below $0.5$ because the tuned encoders
produce more confidently-calibrated probabilities, so a relatively low
decision boundary captures nearly all phishing examples without
inflating false positives.

\subsection{Metrics and their fairness}
We report a broad panel of metrics, appropriate for an imbalanced,
security-critical task:
\begin{itemize}
  \item \textbf{Accuracy} --- an obvious but misleading metric
  under imbalance; reported for completeness only.
  \item \textbf{Precision, Recall, F1-score} --- precision measures
  how many flagged pages are actually phishing, recall measures
  how many phishing pages we catch, and F1 is their harmonic mean.
  \item \textbf{ROC-AUC} --- a threshold-free ranking metric.
  \item \textbf{PR-AUC} --- the area under the precision--recall
  curve; preferable to ROC-AUC on imbalanced data because it
  directly rewards precision on the positive class.
  \item \textbf{C-index} --- concordance-index over probability
  scores, which measures how often a random phishing page is
  scored above a random legitimate page.
  \item \textbf{Confusion matrix} --- exposes the raw counts of
  TP, FP, FN, and TN.
\end{itemize}

\noindent\textbf{Which metric should we focus on?}
Given our problem, we argue that \emph{recall} and \emph{F1 under a
recall-leaning threshold} are the most operationally meaningful
metrics. The cost of a false negative (user credentials stolen)
vastly exceeds the cost of a false positive (user must click past a
warning). ROC-AUC and PR-AUC are useful secondary metrics because
they characterise the quality of the probability ranking
independent of the chosen threshold. Accuracy, by contrast, is
almost uninformative under imbalance: the $2{:}1$ class ratio alone
allows a trivial baseline to reach $65\%$ accuracy. We therefore
report accuracy for completeness but explicitly do not use it as a
model-selection criterion.

\subsection{Baselines}
We train four single-modality baselines that reuse the same
preprocessing, training loop, class weighting, and threshold
calibration: URL-only (BiLSTM + 9 scalar features), Text-only
(frozen DistilBERT + MLP head), Visual-only (frozen EfficientNet-B0
+ MLP head), and HTML-only (8 structural features + MLP head).

\subsection{Ablation and fusion-strategy studies}
We run leave-one-out ablations by setting
\texttt{disabled\_modalities} in the config to each of
\texttt{url}, \texttt{text}, \texttt{visual}, and \texttt{html} in
turn, retraining each variant from scratch. A separate sub-study
compares the three fusion strategies (concatenation, weighted,
attention) under otherwise identical hyperparameters.

\section{Results}
\subsection{Main results: baseline vs.\ Optuna-tuned}
We evaluate every model in two regimes.  The \emph{baseline} regime
uses the defaults in \texttt{configs/config.yaml} and runs a single
training trial per model (attention fusion, $d=256$, AdamW at
$5\!\times\!10^{-4}$, batch size $16$, early stopping with patience
$10$).  The \emph{Optuna} regime runs a $20$-trial TPE search per
unimodal model, promotes the best unimodal hyperparameters into a
second $20$-trial fusion search, and selects the final checkpoint by
the composite score
$s=0.5\cdot\text{F1}+0.25\cdot\text{ROC-AUC}+0.25\cdot\text{C-index}$.
Table~\ref{tab:main} reports test-set performance for both regimes on
the same held-out $246$-sample split.

\begin{table*}[t]
\caption{Test-set performance across the baseline (single-trial,
default config) and Optuna-tuned ($20$-trial TPE per model) regimes.
All numbers are computed on the same held-out $246$-sample split.
Best value per column is bold.}
\label{tab:main}
\centering
\small
\begin{tabular}{llccccccc}
\toprule
Regime & Model & Accuracy & Precision & Recall & F1 & ROC-AUC & PR-AUC & Threshold \\
\midrule
\multirow{5}{*}{Baseline}
 & URL-only (BiLSTM + 9 feats) & 0.841 & 0.724 & 0.884 & 0.796 & 0.921 & 0.885 & 0.36 \\
 & Text-only (DistilBERT)      & 0.764 & 0.606 & \textbf{0.930} & 0.734 & 0.866 & 0.760 & 0.46 \\
 & Visual-only (EfficientNet)  & 0.711 & 0.576 & 0.663 & 0.616 & 0.760 & 0.609 & 0.59 \\
 & HTML-only (8 features)      & 0.646 & 0.495 & 0.628 & 0.554 & 0.707 & 0.552 & 0.58 \\
 & Multi-modal (attention)     & 0.874 & 0.867 & 0.756 & 0.807 & 0.934 & 0.895 & 0.61 \\
\midrule
\multirow{5}{*}{Optuna-tuned}
 & URL-only (CNN + 9 feats)    & 0.886 & 0.882 & 0.779 & 0.827 & 0.940 & 0.920 & 0.28 \\
 & Text-only (DistilBERT)      & 0.813 & 0.744 & 0.709 & 0.726 & 0.875 & 0.754 & 0.42 \\
 & Visual-only (EfficientNet)  & 0.679 & 0.541 & 0.535 & 0.538 & 0.738 & 0.542 & 0.58 \\
 & HTML-only (8 features)      & 0.614 & 0.469 & 0.779 & 0.585 & 0.706 & 0.543 & 0.46 \\
 & \textbf{Multi-modal (weighted)} & \textbf{0.927} & \textbf{0.905} & 0.884 & \textbf{0.894} & \textbf{0.957} & \textbf{0.946} & 0.43 \\
\bottomrule
\end{tabular}
\end{table*}

Three observations stand out.  First, within each regime the fusion
model beats the strongest unimodal branch.  In the baseline regime
attention fusion delivers F1 $=0.807$ vs.\ URL-only's $0.796$
($+0.011$), lifts precision from $0.724$ to $0.867$, and cuts false
positives from $29$ down to $10$.  Second, Optuna tuning is
\emph{strictly} better than the baseline for fusion on every metric:
F1 jumps by $+0.087$ ($0.807\!\to\!0.894$), ROC-AUC by $+0.023$
($0.934\!\to\!0.957$), and PR-AUC by $+0.051$
($0.895\!\to\!0.946$).  The tuned fusion commits $8$ false positives
and $10$ false negatives out of $246$ samples --- a near-balanced and
very low error profile.  Third, the Optuna search did not merely
shuffle learning rates: it selected a \emph{different} URL encoder
(CNN rather than BiLSTM), a smaller $32$-dim character embedding,
one-cycle learning-rate scheduling, weighted random sampling, and
--- crucially --- the \emph{weighted-sum} fusion strategy in place of
multi-head attention.  We analyse the fusion-strategy choice in
Section~VII-C.

\begin{figure}[t]
\centering
\includegraphics[width=0.95\columnwidth]{outputs/evaluation/roc_curves.png}
\caption{ROC curves for all ten evaluated runs (baseline and
Optuna-tuned).  The tuned fusion curve dominates (AUC $=0.957$),
followed by the baseline fusion ($0.934$) and the tuned URL-only
($0.940$).}
\label{fig:roc}
\end{figure}

\begin{figure}[t]
\centering
\includegraphics[width=0.95\columnwidth]{outputs/evaluation/multimodal_(optimization)_cm.png}
\caption{Confusion matrix for the Optuna-tuned multi-modal weighted
fusion on the $246$-sample test set at $T^{*}=0.43$.  Only $8$
false positives and $10$ false negatives.}
\label{fig:cm}
\end{figure}

\subsection{Optuna-tuned hyperparameters}
The promoted URL-branch configuration selected by the unimodal search
(Trial $13$, composite $s=0.907$) replaces the default BiLSTM with a
$1$-D character CNN, halves the character embedding from $64$ to
$32$, widens the URL hidden dimension to $256$, sets dropout to
$0.223$, and retains the nine hand-crafted scalar features (which
the ablation in Section~VII-C confirms are mildly redundant but not
harmful).  The fusion search (Trial $13$, composite $s=0.909$)
picks the weighted-sum fusion strategy with fusion dropout $0.470$,
cosine learning-rate schedule, weighted random sampling, learning
rate $2.0\!\times\!10^{-3}$, and weight decay $1.1\!\times\!10^{-4}$
--- all of which differ from the defaults.  Table~\ref{tab:best_params}
summarises the selected settings.

\begin{table}[t]
\caption{Selected hyperparameters for the tuned multi-modal model
(best Optuna trial, composite $s=0.909$).}
\label{tab:best_params}
\centering
\begin{tabular}{ll}
\toprule
Hyperparameter & Value \\
\midrule
URL encoder architecture   & $1$-D CNN \\
URL char embedding dim     & $32$ \\
URL hidden dim             & $256$ \\
URL dropout                & $0.223$ \\
Use URL scalar features    & \texttt{true} \\
Fusion strategy            & weighted sum \\
Fusion projected dim       & $256$ \\
Fusion hidden dim          & $768$ \\
Fusion dropout             & $0.470$ \\
Optimiser / scheduler      & AdamW / cosine \\
Learning rate              & $2.00\!\times\!10^{-3}$ \\
Weight decay               & $1.12\!\times\!10^{-4}$ \\
Batch size                 & $16$ \\
Sampling strategy          & weighted \\
\bottomrule
\end{tabular}
\end{table}

\subsection{Ablation on the tuned fusion}
With the tuned fusion checkpoint in hand we run a formal leave-one-out
ablation: for each variant we disable a single modality (or swap the
fusion strategy) and re-train from scratch under otherwise identical
tuned hyperparameters.  Table~\ref{tab:ablation} reports the
test-set metrics and the change in composite score relative to the
full tuned model.  The full tuned configuration retrains to F1
$=0.849$, ROC-AUC $=0.954$, PR-AUC $=0.940$, accuracy $=0.890$ at
threshold $T^{*}=0.31$.  Differences from the headline Table~\ref{tab:main}
numbers are due to (i) a fresh training seed and (ii) the fact that
the ablation harness re-calibrates thresholds on a slightly different
training schedule; the overall picture is robust.

\begin{table*}[t]
\caption{Leave-one-out modality ablation and fusion-strategy comparison
on the Optuna-tuned fusion.  $\Delta s$ is the change in composite score
(relative to the \emph{full tuned} run) on the test split.  A negative
$\Delta s$ means the variant is \emph{worse} than the full tuned model.}
\label{tab:ablation}
\centering
\small
\begin{tabular}{lccccccr}
\toprule
Variant & Active modalities & Strategy & F1 & ROC-AUC & PR-AUC & Accuracy & $\Delta s$ vs.\ full \\
\midrule
full\_tuned         & url,text,visual,html & weighted & 0.849 & 0.954 & 0.940 & 0.890 & $\phantom{-}0.000$ \\
no\_url             & text,visual,html     & weighted & 0.709 & 0.889 & 0.781 & 0.797 & $-0.103$ \\
no\_text            & url,visual,html      & weighted & 0.854 & 0.943 & 0.935 & 0.898 & $-0.003$ \\
no\_visual          & url,text,html        & weighted & 0.881 & 0.956 & 0.944 & 0.923 & $+0.017$ \\
no\_html            & url,text,visual      & weighted & 0.826 & 0.947 & 0.929 & 0.878 & $-0.015$ \\
no\_url\_scalar\_features & all (scalars off) & weighted & 0.874 & 0.953 & 0.936 & 0.911 & $+0.012$ \\
strategy\_attention & all                  & attention & 0.875 & 0.940 & 0.939 & 0.911 & $+0.006$ \\
strategy\_weighted  & all                  & weighted  & 0.876 & 0.950 & 0.935 & 0.915 & $+0.011$ \\
strategy\_concatenation & all              & concat    & 0.868 & 0.967 & 0.954 & 0.902 & $+0.016$ \\
\bottomrule
\end{tabular}
\end{table*}

Three findings from the ablation deserve emphasis.  \textbf{(i)~URL
is indispensable.} Removing the URL branch collapses composite score
by $0.103$ and F1 by $0.140$: the no\_url variant falls to F1
$=0.709$, ROC-AUC $=0.889$ --- worse than every individual unimodal
branch except HTML.  This is by far the largest ablation effect in
the study and directly answers RQ4.  \textbf{(ii)~Visual is
redundant on this corpus.}  Removing the visual branch
\emph{improves} composite score by $0.017$ and lifts F1 to $0.881$.
This is consistent with the visual-only baseline being the weakest
unimodal model (F1 $=0.538$ tuned, $0.616$ baseline); on our
$1{,}608$-sample matched corpus the frozen EfficientNet-B0
embeddings contribute more noise than signal, and the fusion stack
slightly benefits from being denied that information.  \textbf{(iii)
Fusion strategy matters less than which modalities are fused.}  All
three fusion strategies (weighted, attention, concatenation) come
within $0.011$ composite score of each other; concatenation
actually gives the highest ROC-AUC ($0.967$) and PR-AUC ($0.954$)
in the retrain, but weighted sum was selected by Optuna because it
maximises the composite score under the tuning objective.  Scalar
URL features (\texttt{use\_url\_scalar\_features=false}) turn out to
be slightly redundant with the learned CNN encoder, consistent with
the CNN having enough capacity to re-derive basic lexical statistics.

\begin{figure}[t]
\centering
\includegraphics[width=0.95\columnwidth]{outputs/optimization/ablation/ablation_metrics.png}
\caption{Ablation metrics (test F1, ROC-AUC, PR-AUC, accuracy) for
each variant in Table~\ref{tab:ablation}.  no\_url is the clear
negative outlier; no\_visual and strategy\_concatenation are slight
positive outliers.}
\label{fig:ablation}
\end{figure}

\subsection{Explainability: SHAP and modality Shapley}
\label{sec:explain}
Beyond ablation, we run a two-level interpretability analysis on the
tuned fusion checkpoint.  A \emph{score-level SHAP} explainer fits an
interpretable $26$-feature surrogate to the fusion model's phishing
probability, reaching $R^{2}=0.955$ and MAE $=0.046$ on the $246$
test samples; this confirms the surrogate is faithful enough to
reason about.  At the feature level the top six contributors by mean
absolute SHAP value are \texttt{url\_phishing\_prob} ($0.217$),
\texttt{max\_unimodal\_prob} ($0.097$), \texttt{visual\_phishing\_prob}
($0.046$), \texttt{confidence} ($0.036$),
\texttt{num\_modalities\_above\_threshold} ($0.031$), and
\texttt{text\_phishing\_prob} ($0.017$).  The URL branch's own
phishing probability is more than twice as influential as the next
strongest feature --- it is, essentially, the single biggest driver
of the fusion output.

A second, architecture-aware analysis computes \emph{modality
Shapley values} directly on the projected $256$-dim fusion vectors
by enumerating all $2^{4}=16$ subsets of
$\{\text{url},\text{text},\text{visual},\text{html}\}$ and measuring
each modality's marginal contribution to the final phishing score.
Averaged over the test set, the mean absolute modality Shapley values
are: URL $=0.158$, visual $=0.094$, text $=0.059$, HTML $=0.036$
(Table~\ref{tab:shapley}).  Two cross-consistency observations are
worth highlighting: (a)~URL is the dominant contributor under
\emph{both} the ablation (Table~\ref{tab:ablation}) and the Shapley
analysis, which strongly corroborates RQ1; and (b)~visual scores
second on Shapley mean-absolute contribution even though removing it
improves composite score (ablation) --- meaning visual contributes
high-magnitude but low-quality evidence that the fusion classifier
has to actively correct for.  This is a textbook case where
importance (how much a feature is used) and usefulness (whether using
it helps) diverge, and it is exactly the kind of finding that
justifies running both analyses.

\begin{table}[t]
\caption{Modality Shapley values on the projected fusion vectors
(tuned model, $246$-sample test set).}
\label{tab:shapley}
\centering
\begin{tabular}{lcc}
\toprule
Modality & Mean Shapley value & Mean $|\Phi|$ \\
\midrule
URL    & $+0.085$ & $0.158$ \\
Visual & $+0.068$ & $0.094$ \\
Text   & $+0.010$ & $0.059$ \\
HTML   & $+0.002$ & $0.036$ \\
\bottomrule
\end{tabular}
\end{table}

\begin{figure}[t]
\centering
\includegraphics[width=0.95\columnwidth]{outputs/explainability/global_importance.png}
\caption{Global SHAP importance for the tuned fusion surrogate
($R^{2}=0.955$).  \texttt{url\_phishing\_prob} dominates the
ranking; HTML-structural and URL-lexical scalars occupy the long
low-importance tail.}
\label{fig:shap}
\end{figure}

\begin{figure}[t]
\centering
\includegraphics[width=0.95\columnwidth]{outputs/explainability/fusion_modality_contributions.png}
\caption{Mean absolute modality Shapley value on projected fusion
vectors.  URL $=0.158$, visual $=0.094$, text $=0.059$,
HTML $=0.036$.}
\label{fig:modality_shap}
\end{figure}

\subsection{Unimodal ranking}
The unimodal half of Table~\ref{tab:main} already answers the
``which modality is strongest?'' question directly, in both regimes.
In the baseline regime the ranking by F1 is
URL ($0.796$) $>$ Text ($0.734$) $>$ Visual ($0.616$) $>$ HTML
($0.554$).  After Optuna tuning the gap widens significantly: URL
lifts to F1 $=0.827$ (and ROC-AUC $=0.940$), while text barely moves
($0.726$), visual actually drops to $0.538$ (the search picked a
lower-dropout, lower-threshold configuration that reduced
recall), and HTML lifts to $0.585$ driven by higher recall
($0.779$).  URL is the only modality that \emph{convincingly}
improves under tuning, confirming that it is both the strongest and
the most tunable signal on this corpus --- a result that the
modality-Shapley analysis in Section~VII-D corroborates.

\subsection{Temperature calibration}
After applying temperature scaling the Brier score on the test set
drops and the prediction-confidence histogram
(\texttt{outputs/evaluation/prediction\_confidence.png}) shifts
visibly away from the two extreme bins. In particular, the number
of confidently-wrong predictions (confidence $>0.9$ but incorrect)
decreases after calibration, which is precisely the behaviour
expected from a well-calibrated classifier and which is important
for the threshold sweep to be meaningful.

\subsection{Threshold sweep}
Fig.~\ref{fig:threshold} plots F1, precision, and recall as a
function of the decision threshold on the validation set for the
baseline fusion model.  The precision curve is monotonically
non-decreasing in the threshold, and the recall curve is
monotonically non-increasing.  F1 peaks between $0.55$ and $0.7$ for
the baseline multi-modal model, giving $T^{*}=0.61$ in that regime.
The tuned fusion model, whose probabilities are better calibrated
(temperature $=1.55$), selects a markedly lower threshold of
$T^{*}=0.43$ --- not because the classifier is less confident, but
because the tuned model's probability distribution is narrower and
more symmetric around the decision boundary, so the F1-optimal
cutoff sits closer to $0.5$.

\begin{figure}[t]
\centering
\includegraphics[width=0.95\columnwidth]{outputs/evaluation/threshold_sweep.png}
\caption{Precision--recall--F1 vs.\ decision threshold on the
multi-modal model.}
\label{fig:threshold}
\end{figure}

\subsection{Misclassification analysis}
Inspecting the $8$ false positives and $10$ false negatives of the
\emph{tuned} multi-modal model is revealing.  The top false
positives are cryptic-looking but benign URLs containing tildes,
numeric IDs, query strings, and deep paths (e.g.,
\texttt{www.macs.hw.ac.uk/\textasciitilde pdw/},
\texttt{arasgallery.com/profile.php?id=40}), which trigger the
same lexical cues the model associates with phishing.  Several
top false negatives are short, innocent-looking phishing URLs
hosted on URL shorteners or low-reputation hosts (e.g.,
\texttt{tiny.cc/cartasi-it}, \texttt{tide.org.au/quick/}).  These
are genuinely hard cases: the URL alone is insufficient, and
without a screenshot that imitates a recognised brand the visual
branch also struggles.  Both error modes suggest that
host-reputation and DNS-age features, which we do not currently
use, would be fruitful extensions.  A full list of the top
errors is saved to
\texttt{outputs/evaluation/misclassification\_analysis.txt}.

\section{Results Discussion}
\subsection{How the loss evolved over time}
Fig.~\ref{fig:curves} summarises training and validation dynamics
for the multi-modal model. Training loss decreases smoothly from
approximately $0.68$ at epoch 1 to below $0.30$ by epoch 17, while
validation loss decreases more erratically --- dropping to a
minimum in the first few epochs and then oscillating around
$0.5$--$0.6$. Validation F1 climbs from around $0.58$ at epoch 1
to a plateau in the $0.75$--$0.82$ range, and early stopping
halts training at epoch $17$.

\begin{figure}[t]
\centering
\includegraphics[width=0.95\columnwidth]{outputs/multimodal/training_curves_multimodal_fusion.png}
\caption{Multi-modal training curves. Training loss decreases
smoothly; validation F1 climbs into the $0.75$--$0.82$ range.
Early stopping halts training at epoch $17$.}
\label{fig:curves}
\end{figure}

Two observations are worth highlighting. First, the gap between
training and validation loss widens after the first few epochs,
which is the classic signature of mild over-fitting. This is
expected given our small dataset and motivates the dropout ($0.3$),
LayerNorm, and augmentation we employ. Second, the checkpoint
selected by the early-stopping callback is \emph{not} the final
epoch --- it is the best-validation-loss checkpoint. Using the best
checkpoint rather than the last-epoch model is a simple but
important implementation choice that cleanly separates model
selection from test-time evaluation.

\subsection{Did the model actually learn something meaningful?}
We check this by visualising the text+visual embedding space using
$t$-SNE (Fig.~\ref{fig:tsne}). The two classes do \emph{not} form
perfectly separated clusters --- a substantial mixed region exists,
which matches the intuition that raw pre-trained embeddings carry
only partial phishing signal. However, several regions of the
projection are clearly dominated by one class (note the
phishing-heavy region in the lower left), and the mixed region
corresponds almost exactly to the hard examples identified in the
misclassification analysis. This qualitative result is consistent
with the quantitative finding that text and visual alone produce
the weakest unimodal classifiers; it is only after combining these
signals with the URL branch and passing them through the
attention-fusion module that the fused representation becomes
cleanly separable, which is why the fusion model outperforms every
unimodal baseline.

\subsection{What the attention weights tell us (baseline fusion)}
The \emph{baseline} fusion uses the attention strategy and therefore
exposes a $4\times 4$ attention-weight matrix for every example,
which we average over the test set and plot in
Fig.~\ref{fig:attention}.  The dominant pattern is that \emph{Text
is the most attended-to modality overall}: the Visual query attends
to Text with weight $0.452$, URL attends to Text with $0.355$, and
HTML attends to Visual with $0.382$.  HTML tokens receive the
lowest average attention across every query.  This is a more
nuanced picture than ``URL dominates'': Text acts as the semantic
hub that other modalities consult, while URL is strongest as a
standalone predictor.  The Optuna search swapped this mechanism for
weighted-sum fusion, where the learned softmax weights over the four
modalities converged to a clear URL-dominant pattern --- which is
precisely what the modality-Shapley analysis in Section~VII-D
quantifies without needing attention weights to interpret.

\begin{figure}[t]
\centering
\includegraphics[width=0.95\columnwidth]{outputs/evaluation/attention_weights.png}
\caption{Average attention weight heatmap across the test set for
the \emph{baseline} attention fusion.  Rows are query modalities,
columns are key modalities.  Text is the most attended-to modality
(highest column weights); HTML is the least attended-to.}
\label{fig:attention}
\end{figure}

\subsection{Prediction-confidence histogram}
The prediction-confidence histogram (Fig.~\ref{fig:confidence}) shows
a strongly bimodal distribution, with most legitimate predictions
clustered near $0.1$ and most phishing predictions clustered near
$1.0$. The ``uncertain'' mass in the $[0.3, 0.6]$ range is small
but concentrated around the hardest test examples. A well-shaped
confidence histogram is desirable because it indicates that the
classifier's output is a meaningful probability rather than an
arbitrary logit --- which in turn is what makes threshold-based
policies such as ``warn the user above $T^{*}=0.43$'' operationally
trustworthy.

\begin{figure}[t]
\centering
\includegraphics[width=0.95\columnwidth]{outputs/evaluation/prediction_confidence.png}
\caption{Prediction confidence distribution for the multi-modal
model. Bimodal shape indicates well-separated classes.}
\label{fig:confidence}
\end{figure}

\subsection{Feature correlation}
The feature-correlation heatmap (Fig.~\ref{fig:featurecorr}) shows
the Pearson correlation between each hand-crafted feature and the
label. The strongest positive correlates are
\texttt{slashes} ($+0.26$), \texttt{dots} ($+0.23$), and
\texttt{ext\_link\_ratio} ($+0.21$); the strongest negative correlates
are \texttt{text\_length} and \texttt{has\_password} (both $-0.26$).
\texttt{has\_https} has effectively zero variance in this corpus
(shown as a blank row in the heatmap), confirming that TLS alone is
no longer a reliable indicator. All individual correlations are
modest ($|r|<0.27$), which is precisely why no single-feature or
single-modality classifier is enough and why a learned non-linear
fusion over multiple weak signals --- the very structure of our
architecture --- is the right approach.

\subsection{Answers to the research questions}
We can now answer the four research questions posed in Section II.

\noindent\textbf{RQ1 --- Which modality is strongest?}
URL, in every evaluation we ran.  It is the strongest unimodal
branch in both regimes (baseline F1 $=0.796$, tuned F1 $=0.827$),
it has the largest leave-one-out drop in the ablation (composite
$\Delta s = -0.103$, Table~\ref{tab:ablation}), the largest mean
absolute SHAP contribution at the score level
($|$SHAP$|=0.217$, Fig.~\ref{fig:shap}), and the largest mean
absolute modality Shapley value at the representation level
(Table~\ref{tab:shapley}, $|\Phi|=0.158$).  Four independent
analyses converging on the same answer is the kind of
triangulation that makes us confident in it.

\noindent\textbf{RQ2 --- Does multi-modal fusion help?}
Yes, and tuning amplifies the gap.  In the baseline regime attention
fusion already beats URL-only on F1 ($0.807$ vs.\ $0.796$),
precision ($0.867$ vs.\ $0.724$), and PR-AUC ($0.895$ vs.\ $0.885$).
After Optuna tuning the weighted fusion model beats the tuned URL-only
model on \emph{every} headline metric: F1 ($0.894$ vs.\ $0.827$),
ROC-AUC ($0.957$ vs.\ $0.940$), accuracy ($0.927$ vs.\ $0.886$),
precision ($0.905$ vs.\ $0.882$), recall ($0.884$ vs.\ $0.779$),
and PR-AUC ($0.946$ vs.\ $0.920$).  The tuned fusion commits $8$
false positives and $10$ false negatives on $246$ samples --- a
deployment-relevant operating point where both error types are
low.

\begin{figure}[t]
\centering
\includegraphics[width=0.95\columnwidth]{outputs/evaluation/pr_curves.png}
\caption{Precision--recall curves for all baseline and tuned
models.  The tuned fusion curve dominates (PR-AUC $=0.946$).}
\label{fig:pr}
\end{figure}

\noindent\textbf{RQ3 --- Which fusion strategy is best?}
The Optuna search picked \emph{weighted sum} (strategy = weighted,
$d=256$, hidden $=768$, dropout $=0.470$), and the retrained
ablation confirms that weighted sum edges out attention under the
composite metric ($\Delta s = +0.011$ vs.\ $+0.006$).  However, all
three strategies land within $0.011$ composite score of each other;
concatenation actually gives the highest ROC-AUC ($0.967$) in the
retrain.  The practical conclusion is that \emph{which} modalities
are fused matters much more than \emph{how} they are fused: the URL
branch drives the result, and any of the three fusion mechanisms
is capable of exploiting it.

\noindent\textbf{RQ4 --- How much does each modality contribute?}
The formal ablation gives a clean quantitative answer
(Table~\ref{tab:ablation}): removing URL costs composite score
$-0.103$ (F1 drops to $0.709$), removing HTML costs $-0.015$,
removing text costs $-0.003$, and removing visual
\emph{improves} composite score by $+0.017$.  The modality-Shapley
analysis (Table~\ref{tab:shapley}) agrees on URL being dominant
but ranks the other modalities by absolute magnitude rather than by
``usefulness'': visual has the second-largest $|\Phi|$ despite
being the one ablation variant whose removal helps.  This is
consistent with an interpretation in which visual pushes the fusion
score around more than text or HTML do on average, but not
productively --- the classifier has to learn to down-weight it.

\begin{figure}[t]
\centering
\includegraphics[width=0.95\columnwidth]{outputs/evaluation/metrics_comparison.png}
\caption{Side-by-side comparison of key metrics across all models.}
\label{fig:metrics}
\end{figure}

\subsection{Computational trade-offs}
The multi-modal model has $1.06$M trainable parameters and trains
in under 10 minutes on CPU for our dataset, with per-epoch times
of roughly 20--25 seconds. Inference is dominated by the frozen
DistilBERT forward pass on the webpage text; URL, HTML, and
visual forward passes are comparatively negligible. For
latency-sensitive deployments it is straightforward to pre-cache
text and visual embeddings at crawl time, which is why the
project also implements a \texttt{FastFusionClassifier} variant
that consumes pre-extracted DistilBERT and EfficientNet
embeddings directly.

\subsection{Limitations and threats to validity}
Our dataset has three important limitations. First, 1{,}608
samples is modest and only 246 are in the held-out test set, so
the confidence intervals on our reported metrics are wide. Second,
both source datasets are slightly biased: the HuggingFace URL
corpus skews towards older phishing sites and the Kaggle
screenshot corpus is dominated by a single browser configuration.
Third, the class balance is $\approx 2{:}1$, which is gentler
than the distribution a real deployment would see. Beyond dataset
issues, our frozen-backbone design trades predictive ceiling for
training speed; light fine-tuning of the top layers would almost
certainly improve F1 further.

\section{Conclusion}
We presented a multi-modal phishing detection system that jointly
leverages URL character sequences and hand-crafted lexical
features, webpage text encoded by a frozen DistilBERT, screenshots
encoded by a frozen EfficientNet-B0, and eight hand-crafted HTML
structural features. Each modality is projected into a common
$256$-dim space and combined through a four-head attention fusion
module before a three-layer bottleneck classifier. Beyond model
design, we treated the problem as a careful statistical-learning
exercise: we performed an exploratory data study, addressed class
imbalance with class weighting and threshold calibration, addressed
over-fitting with frozen backbones, dropout, and augmentation, and
tuned hyperparameters.

On a dataset of 1{,}608 matched webpage records the
\emph{baseline} attention-fusion multi-modal model already
out-performs every unimodal branch (F1 $=0.807$, ROC-AUC $=0.934$),
but a full Optuna hyperparameter search lifts the tuned
weighted-sum fusion to F1 $=0.894$, ROC-AUC $=0.957$,
PR-AUC $=0.946$, and accuracy $=0.927$ --- a $+0.087$ F1 gain
over the baseline fusion and a $+0.067$ F1 gain over the tuned
URL-only model.  The tuned fusion also attains the highest
precision ($0.905$) and the fewest false positives ($8$), which
is the most deployment-relevant operating point for a
security-critical application.  A formal leave-one-out ablation
and a two-level SHAP/Shapley analysis both identify URL as the
dominant modality; visual, despite ranking second by mean
absolute Shapley magnitude, is the only modality whose removal
actually improves composite score, which is a finding we would
not have uncovered without running both analyses.  The
multi-modal advantage only emerged cleanly after we corrected
three subtle asymmetries between the unimodal and multimodal
branches --- matched classifier capacity, large-enough projection
dimension, and removing a mixup regulariser that was applied to
only one branch --- a methodological lesson worth highlighting.

Several directions are worth pursuing. First, light fine-tuning of
the top layers of DistilBERT and EfficientNet on phishing data
would likely lift performance further. Second, richer URL and
host-level features (WHOIS age, TLS certificate validity, DNS
reputation) would directly address the false-negative patterns we
observed on URL-shortener and low-reputation phishing pages.
Third, an adversarial robustness study --- perturbing text,
screenshots, or URLs to mimic evasion attempts --- would give a
clearer picture of how each modality behaves under attack.
Finally, deploying the model behind a threshold policy that adapts
to user risk profiles would turn the recall-over-precision
preference we implemented here into a principled user-facing
control.

\section*{Acknowledgements}
We thank the DSA4266 teaching team for their guidance throughout
the project, and the maintainers of PyTorch, HuggingFace
Transformers, torchvision, and timm without which this work would
not have been possible.

\begin{thebibliography}{00}
\bibitem{apwg} Anti-Phishing Working Group, ``Phishing Activity
Trends Reports,'' 2023.
\bibitem{urlnet} H. Le, Q. Pham, D. Sahoo, and S. C. H. Hoi,
``URLNet: Learning a URL Representation with Deep Learning for
Malicious URL Detection,'' arXiv:1802.03162, 2018.
\bibitem{distilbert} V. Sanh, L. Debut, J. Chaumond, and T. Wolf,
``DistilBERT, a distilled version of BERT: smaller, faster, cheaper
and lighter,'' in \emph{NeurIPS EMC$^2$ Workshop}, 2019.
\bibitem{effnet} M. Tan and Q. V. Le, ``EfficientNet: Rethinking
Model Scaling for Convolutional Neural Networks,'' in \emph{ICML},
2019.
\bibitem{bert} J. Devlin, M.-W. Chang, K. Lee, and K. Toutanova,
``BERT: Pre-training of Deep Bidirectional Transformers for
Language Understanding,'' in \emph{NAACL-HLT}, 2019.
\bibitem{resnet} K. He, X. Zhang, S. Ren, and J. Sun, ``Deep
Residual Learning for Image Recognition,'' in \emph{CVPR}, 2016.
\bibitem{attn} A. Vaswani \textit{et al.}, ``Attention Is All You
Need,'' in \emph{NeurIPS}, 2017.
\bibitem{phishrnn} A. C. Bahnsen \textit{et al.}, ``Classifying
phishing URLs using recurrent neural networks,'' in \emph{eCrime},
2017.
\bibitem{visualphish} S. Abdelnabi, K. Krombholz, and M. Fritz,
``VisualPhishNet: Zero-Day Phishing Website Detection by Visual
Similarity,'' in \emph{ACM CCS}, 2020.
\bibitem{multimodal-survey} T. Baltru\v{s}aitis, C. Ahuja, and
L.-P. Morency, ``Multimodal Machine Learning: A Survey and
Taxonomy,'' \emph{IEEE TPAMI}, 2019.
\bibitem{mixup} H. Zhang, M. Cisse, Y. N. Dauphin, and D.
Lopez-Paz, ``mixup: Beyond Empirical Risk Minimization,'' in
\emph{ICLR}, 2018.
\bibitem{tempscale} C. Guo, G. Pleiss, Y. Sun, and K. Q.
Weinberger, ``On Calibration of Modern Neural Networks,'' in
\emph{ICML}, 2017.
\bibitem{optuna} T. Akiba \textit{et al.}, ``Optuna: A Next-
generation Hyperparameter Optimization Framework,'' in
\emph{KDD}, 2019.
\end{thebibliography}

\end{document}
